\section{Passive cut set and gravitational endpoint resource}
\label{sec:networkmaterial}

\subsection{Passive selected-port cut}

Use an energy-normalized passive modal realization
\begin{equation}
\dot{\bm a}=A\bm a-K^\dagger\bm b_{\rm in},
\qquad
\bm b_{\rm out}=\bm b_{\rm in}+K\bm a,
\label{eq:passiveio}
\end{equation}
with
\begin{equation}
\boxed{A=-iH-\frac12K^\dagger K,\qquad H=H^\dagger.}
\label{eq:Apassive}
\end{equation}
Partition $K=(K_u,K_g,K_\ell)^T$ into useful, gravitational, and other passive ports. The useful-input to gravitational-output block is
\begin{equation}
S_{g\leftarrow u}(s)=-K_g(sI-A)^{-1}K_u^\dagger,
\label{eq:Sgu}
\end{equation}
and for a strictly proper transfer define
\begin{equation}
\|S\|_2^2
=\frac1{2\pi}\int_{-\infty}^{\infty}\dd\omega\,
\Tr[S^\dagger(i\omega)S(i\omega)].
\label{eq:H2def}
\end{equation}
Passivity gives the selected controllability Gramian $0\le P_u\le I$, hence \cite{GutaYamamoto2014,TechakesariNurdin2017}
\begin{equation}
\boxed{
\|S_{g\leftarrow u}\|_2^2
=\Tr(K_gP_uK_g^\dagger)
\le\Tr(K_g^\dagger K_g).}
\label{eq:endpointbound}
\end{equation}
The receiver obeys the dual bound.

The same statement holds directly on a separable modal Hilbert space when the full port operator is bounded and $K_g$ is Hilbert--Schmidt. If $\mathcal T(t)=e^{At}$ is the contraction semigroup, then for every $\tau$
\begin{equation}
0\le P_u(\tau)
\le I-\mathcal T(\tau)\mathcal T^\dagger(\tau)
\le I.
\label{eq:infgramian}
\end{equation}
The monotone strong limit and operator-valued Plancherel recover Eq.~\eqref{eq:endpointbound}. This is standard infinite-dimensional systems machinery \cite{BarasBrockett1975,OpmeerReisWollner2013}; arbitrary unbounded PDE boundary ports require separate admissibility analysis and are outside the present theorem.

Let $P_g(\omega)$ be normalized propagation between the source gravitational output and receiver gravitational input, and define
\begin{equation}
\eta_{\max}
=\sup_{\omega\in\mathcal B}\|P_g(\omega)\|_{\op}^2.
\label{eq:etamax}
\end{equation}
For
\begin{equation}
T(\omega)=S_{v\leftarrow g}^{(B)}(i\omega)
P_g(\omega)S_{g\leftarrow u}^{(A)}(i\omega),
\label{eq:fullT}
\end{equation}
applying the Hilbert--Schmidt/operator-norm inequality at either endpoint gives
\begin{theorem}[Passive cut set]
\begin{equation}
\boxed{
\Gamma_{\rm coh}
\le
\eta_{\max}
\min\!\left[
\Tr(K_{g,A}^\dagger K_{g,A}),
\Tr(K_{g,B}^\dagger K_{g,B})
\right].}
\label{eq:cutset}
\end{equation}
\end{theorem}
This theorem is generic passive systems mathematics; the gravitational content enters through the two resources on its right-hand side.

\subsection{Microscopic port normalization}

Let $G_A:\mathcal H_{m,A}\to\mathcal H_g$ be the narrowband matter-to-radiation coupling operator. Its damping Gram operator is $\Gamma_{g,A}=G_A^\dagger G_A$. Polar decomposition gives
\begin{equation}
G_A=V_A\Gamma_{g,A}^{1/2},
\qquad
G_B=V_B\Gamma_{g,B}^{1/2},
\label{eq:polar}
\end{equation}
so fixed-frequency propagation $U_R$ factors as
\begin{equation}
\boxed{
G_B^\dagger U_RG_A
=\Gamma_{g,B}^{1/2}P_g\Gamma_{g,A}^{1/2},
\qquad P_g=V_B^\dagger U_RV_A.}
\label{eq:microscopicfactor}
\end{equation}
Thus linewidth/oscillator strength and normalized radiation geometry are not counted twice. Nonorthogonal radiation patterns simply appear as off-diagonal entries of $\Gamma_g$.

\subsection{Inertia bound on the total gravitational port}

Let $\rho(\bm r)$ be the mass density, with origin at the center of mass, and use
\begin{equation}
\langle\bm u,\bm v\rangle_\rho
=\int\rho\,\bm u\cdot\bm v\,dV.
\label{eq:massinner}
\end{equation}
For mutually orthogonal normal modes $\bm w_n$ with modal masses $\mu_n$, the linearized STF quadrupole is
\begin{equation}
\delta Q_{ij}[\bm u]
=\int\rho\left(u_ix_j+u_jx_i-\frac23\delta_{ij}\bm u\cdot\bm x\right)dV,
\label{eq:linearQmap}
\end{equation}
with influence fields
\begin{equation}
(g^{ij})_k
=\delta_{ik}x_j+\delta_{jk}x_i-\frac23\delta_{ij}x_k
\label{eq:gijfield}
\end{equation}
and modal quadrupoles
\begin{equation}
q_{n,ij}=\langle\bm w_n,\bm g^{ij}\rangle_\rho.
\label{eq:qmodal}
\end{equation}
These are the standard long-wavelength tidal-mode couplings of resonant-mass antenna theory \cite{Lobo1995}. The STF contraction gives
\begin{equation}
\boxed{
\sum_{ijk}|(g^{ij})_k|^2=\frac{20}{3}r^2,}
\label{eq:gidentity}
\end{equation}
and Bessel completeness therefore yields
\begin{equation}
\boxed{
\sum_n\frac{q_n:q_n}{\mu_n}\le\frac{20}{3}I,
\qquad I=\int\rho r^2dV.}
\label{eq:classicalqsum}
\end{equation}
The method is standard modal-participation/effective-mass completeness \cite{BamfordGaymanWada1971}; the point here is its gravitational STF specialization.

Hirakawa, Narihara, and Fujimoto define the gravitational effective area \cite{HirakawaNariharaFujimoto1976}
\begin{equation}
A_{Gn}=\frac{2q_n:q_n}{M\mu_n},
\label{eq:HirakawaAG}
\end{equation}
so Eq.~\eqref{eq:classicalqsum} becomes
\begin{equation}
\boxed{\sum_nMA_{Gn}\le\frac{40}{3}I.}
\label{eq:classicalAGsum}
\end{equation}
For a harmonic mode, their radiated-power normalization gives the gravitational energy-decay rate
\begin{equation}
\boxed{
\kappa_{g,n}=\frac{GMA_{Gn}\omega_n^4}{10c^5}.}
\label{eq:HirakawaKappa}
\end{equation}
Consequently
\begin{equation}
\boxed{
\Tr(K_g^\dagger K_g)=\sum_n\kappa_{g,n},}
\label{eq:tracewidths}
\end{equation}
and for retained modes with $\omega_n\le\Omega$,
\begin{equation}
\boxed{
\Tr(K_g^\dagger K_g)
\le\frac{4G}{3c^5}I\Omega^4.}
\label{eq:materialbound}
\end{equation}
This finite trace also makes $K_g$ Hilbert--Schmidt in the countably infinite bounded-port extension. Passive mode mixing cannot increase it because $\Tr(G^\dagger G)$ is basis invariant; the participation factors of recent multimode resonant-mass proposals provide a concrete example of this redistribution \cite{TobarPikovskiTobar2025}.

Quantization is not required for the bound but gives an independent normalization check. For a one-quantum transition,
\begin{equation}
\kappa_{g,n}
=\frac{2G\omega_n^5}{5\hbar c^5}Q_{ij}^{0n}Q_{ij}^{n0},
\label{eq:gravitonrate}
\end{equation}
and quantizing the same normal coordinate gives $Q^{01}:Q^{10}=\hbar M A_{Gn}/(4\omega_n)$, reproducing Eq.~\eqref{eq:HirakawaKappa}. The corresponding mass-quadrupole EWSR yields the same cumulative coefficient; Appendix~\ref{app:sumrule} records that equivalent quantum route.

Combining Eqs.~\eqref{eq:cutset} and \eqref{eq:materialbound} gives
\begin{equation}
\boxed{
\Gamma_{\rm coh}
\le
\eta_{\max}\frac{4G\Omega^4}{3c^5}
\min(I_A,I_B).}
\label{eq:materialcutset}
\end{equation}
It remains only to bound the normalized compact TT propagation factor.